% Bab 8: Technical & AI-Assisted Development Plan
\chapter{Technical \& AI-Assisted Development Plan}
\label{bab:technical-ai-assisted-development-plan}
\section{Proposed Technology}

\textit{Contoh rencana teknologi awal; pilihan teknologi boleh direvisi berdasarkan hasil eksplorasi teknis.}

\begin{longtable}{|L{0.333}|L{0.333}|L{0.333}|}
\caption{Proposed Technology}\\
\hline
\cellcolor[HTML]{D9EAF7}\textbf{Component} & \cellcolor[HTML]{D9EAF7}\textbf{Proposed Technology / Tool} & \cellcolor[HTML]{D9EAF7}\textbf{Purpose} \\ \hline
\endfirsthead
\hline
\cellcolor[HTML]{D9EAF7}\textbf{Component} & \cellcolor[HTML]{D9EAF7}\textbf{Proposed Technology / Tool} & \cellcolor[HTML]{D9EAF7}\textbf{Purpose} \\ \hline
\endhead
Application & React Native + Expo (TypeScript), dibangun dengan Expo Development Build / EAS Build & Pengembangan aplikasi mobile Android dan pembuatan APK untuk perangkat uji \\ \hline
Location & expo-location + expo-task-manager & Membaca lokasi akurasi tinggi (\(\pm\)3 detik / \(\ge\)5 m), pelacakan background melalui foreground service saat Mode Jelajah aktif \\ \hline
Geofence Logic & Modul TypeScript internal (haversine, hysteresis, cooldown) & Mendeteksi entry, proximity, dwell, dan exit event serta mengabaikan pembacaan dengan akurasi \textgreater{}30 m \\ \hline
Map/Navigation & react-native-maps dengan tile OpenStreetMap + overlay denah kawasan; Linking (intent Google Maps) & Menampilkan posisi pengguna, marker POI, garis rute ke tenant, dan navigasi menuju kawasan \\ \hline
Compass & expo-location (watchHeadingAsync) & Menentukan arah panah navigasi menuju tenant \\ \hline
Data & JSON lokal (dibundel di APK) + AsyncStorage & Menyimpan data POI dan konten secara offline, serta preferensi, riwayat kunjungan, dan progres trail di perangkat \\ \hline
Audio/Multimedia & expo-audio, expo-speech (TTS), expo-haptics & Memutar narasi audio, panduan suara navigasi, dan getar saat trigger \\ \hline
Notification & expo-notifications (notification category dengan tombol aksi) & Notifikasi entry/proximity, tombol Dengar/Lewati di lock screen, dan pengingat sisa waktu \\ \hline
Search & Fuse.js & Pencarian toleran typo berdasarkan nama tenant, menu, dan kategori \\ \hline
AI Development Tool & Claude, Claude Code, ChatGPT, Antigravity & Membantu perencanaan, coding, debugging, dan penyusunan test case \\ \hline
Other & expo-camera (QR, opsional untuk F10), Jest + React Native Testing Library, Android Emulator (mock location), Git/GitHub & Scan QR Paspor Kuliner, unit test logika, simulasi pergerakan lokasi, dan version control \\ \hline
\end{longtable}

\section{AI-Assisted Development Plan}

\begin{longtable}{|L{0.250}|L{0.174}|L{0.298}|L{0.279}|}
\caption{AI-Assisted Development Plan}\\
\hline
\cellcolor[HTML]{D9EAF7}\textbf{Aktivitas} & \cellcolor[HTML]{D9EAF7}\textbf{AI / Tool} & \cellcolor[HTML]{D9EAF7}\textbf{Penggunaan yang Direncanakan} & \cellcolor[HTML]{D9EAF7}\textbf{Cara Verifikasi} \\ \hline
\endfirsthead
\hline
\cellcolor[HTML]{D9EAF7}\textbf{Aktivitas} & \cellcolor[HTML]{D9EAF7}\textbf{AI / Tool} & \cellcolor[HTML]{D9EAF7}\textbf{Penggunaan yang Direncanakan} & \cellcolor[HTML]{D9EAF7}\textbf{Cara Verifikasi} \\ \hline
\endhead
Requirement/Planning & Claude, ChatGPT & Membantu mengkritisi user flow dan kelengkapan acceptance criteria & Tim membandingkan hasil dengan PRD dan menyetujui perubahan secara manual \\ \hline
UI/Prototype & AI design assistant & Menghasilkan alternatif wireframe awal & Tim memilih dan menyesuaikan berdasarkan user scenario \\ \hline
Coding & Claude, Antigravity & Membantu membuat boilerplate komponen React Native, modul geofence (haversine, hysteresis, cooldown), algoritma nearest neighbor untuk trail planner, dan scoring rekomendasi & Code review oleh anggota tim melalui pull request; kode dijalankan di perangkat uji dan dicocokkan dengan acceptance criteria F01--F10 \\ \hline
Debugging & Claude, ChatGPT & Membantu menganalisis error log (Metro/Logcat), masalah izin lokasi background, notifikasi, dan pemutaran audio & Perbaikan diuji ulang pada emulator dan perangkat fisik; bug dianggap selesai jika skenario yang gagal berhasil dijalankan kembali \\ \hline
Testing & Claude, Antigravity & Membantu menyusun unit test (Jest) untuk logika geofence, filter, rekomendasi, dan trail planner, serta skenario uji lapangan dan usability test dari acceptance criteria & Tim meninjau kesesuaian test case dengan PRD; hasil akhir diverifikasi melalui uji lapangan di Sudirman Street Food dan dicatat dalam laporan pengujian \\ \hline
\end{longtable}

Prinsip: AI may assist development, but the team remains responsible for understanding, verifying, testing, and demonstrating every implemented feature.
